%
% Front matter
%

\begin{frontmatter}

\title{
  T2K 6.5700 $\times10^{20}$-POT Joint 3-Flavour Oscillation Analysis 
}

\vspace*{0.7cm}

\author[RAL]     {Costas Andreopoulos}
\author[IFIC]    {Anselmo Cervera}
\author[RAL,OXF] {Thomas Dealtry}
\author[IFIC]    {Lorena Escudero S\'anchez}
\author[LANCS]   {Nick Grant}
\author[IFIC]    {Michel Sorel}
\author[IFIC]    {and Panagiotis Stamoulis}

\address[RAL]{Particle Physics Dept., Rutherford Appleton Laboratory, STFC \\
              Harwell Science and Innovation Campus\\
              Oxfordshire OX11 0QX, UK}
\address[OXF]{Subdepartment of Particle Physics, University of Oxford \\
              Denys Wilkinson Building, Keble Road\\
              Oxfordshire OX1 3RH, UK}
\address[LANCS]{Physics Dept., Lancaster University \\
              Lancaster, LA1 4YB, UK}
\address[IFIC]{Instituto de Fisica Corpuscular\\
              CSIC and University of Valencia\\
              Valencia, Spain}

%
% Abstract
%

\begin{abstract}

The results of the {\em VALOR}
joint 3-flavour oscillation analysis on the combined Run 1+2+3+4 dataset 
are reported in this technical note.
This dataset corresponds to an integrated JPARC neutrino beam exposure of 
6.5700 $\times10^{20}$ POT.

This analysis is an update of the work presented in Ref. \cite{OurJointAnalysis} based on the 
{\em VALOR}  $\nu_{\mu}$-disappearance analysis presented in Refs. \cite{OurRun1DisappAnalysis}, 
\cite{OurRun12DisappAnalysis}, \cite{OurRun12DisappAnalysis3Flavour} and \cite{OurRun123DisappAnalysis3Flavour},
with improved methods and an enlarged dataset. 
Furthermore, this updated joint analysis also contains the correlations between the SK efficiencies for the
single $\mu$-like ring and the single e-like ring event samples that were not available in the previous analysis \cite{OurJointAnalysis}.

In this analysis, 
the atmospheric squared-mass splitting $|\Delta m^{2}|$ (the Fogli-Lisi convention is used, see main text for details),
the mixing angles $sin^{2}\theta_{23}$, $sin^{2}\theta_{13}$ 
and the CP-invariance violating phase $\delta_{CP}$,
are determined jointly through a simultaneous oscillation fit to the reconstructed
energy spectrum of the SuperK single $\mu$-like ring and single e-like ring event samples.
The oscillation fit was performed in a 3-flavour framework
including matter effects in constant density matter.
Official T2K 2013 analysis inputs were used.
All 64 systematics considered in this analysis were allowed to float in the fit.


Assuming normal hierachy (NH), 
the joint 3-flavour oscillation fit to
the observed reconstructed energy spectrum of single $\mu$-like ring and single e-like ring
events of the combined Run 1+2+4 dataset,
with all mixing angles constrained to the physical region,
yields
$|\Delta m^{2}|$ = {\color{red}{x.xxx}}$\times$10$^{-3}$ $eV^{2}/c^{4}$,
$\sin^{2}\theta_{23}$ = {\color{red}{x.xxx}},
$\sin^{2}\theta_{13}$ = {\color{red}{x.xxx}} and
$\delta_{CP}$ = {\color{red}{x.xxx}}
as the best-fit values (p-value = {\color{red}{x.xx}}).
The corresponding results for inverted hierachy (IH) are
$|\Delta m^{2}|$ = {\color{red}{x.xxx}} $\times$10$^{-3}$ $eV^{2}/c^{4}$,
$\sin^{2}\theta_{23}$ = {\color{red}{x.xxx}},
$\sin^{2}\theta_{13}$ = {\color{red}{x.xxx}} and
$\delta_{CP}$ = {\color{red}{x.xxx}}
(p-value = {\color{red}{x.xx}}).
Both fits used 2012 PDG values for
$\sin^{2}2\theta_{12}$ and the solar squared-mass splitting $\delta m^{2}$ \cite{PDG}.

Confidence regions are constructed using only the constant-$\Delta\chi^2$ method in the present technical note,
alternative methods, such as Feldman-Cousins \cite{FeldmanCousins}, will be evaluated in following versions.
For NH, the 90\% (68\%) CL allowed region for the four oscillation parameters can be summarized as
{\color{red}{x.xxx}}$\times$10$^{-3}$ ({\color{red}{x.xxx}}$\times$10$^{-3}$) $eV^{2}/c^{4}$
$< |\Delta m^{2}| < $
{\color{red}{x.xxx}}$\times$10$^{-3}$ ({\color{red}{x.xxx}}$\times$10$^{-3}$) $eV^{2}/c^{4}$,
{\color{red}{x.xxx}} ({\color{red}{x.xxx}}) $< sin^{2}\theta_{23} <$ {\color{red}{x.xxx}} ({\color{red}{x.xxx}}),
{\color{red}{x.xxx}} ({\color{red}{x.xxx}}) $< sin^{2}\theta_{13} <$ {\color{red}{x.xxx}} ({\color{red}{x.xxx}}) and
-$\pi$ (-$\pi$) $< \delta_{CP} <$ $\pi$ ($\pi$).
For IH, the corresponding 90\% (68\%) CL allowed regions are
{\color{red}{x.xxx}} $\times$10$^{-3}$ ({\color{red}{x.xxx}}$\times$10$^{-3}$) $eV^{2}/c^{4}$
$< |\Delta m^{2}| < $
{\color{red}{x.xxx}} $\times$10$^{-3}$ ({\color{red}{x.xxx}}$\times$10$^{-3}$) $eV^{2}/c^{4}$,
 {\color{red}{x.xxx}} ({\color{red}{x.xxx}}) $< sin^{2}\theta_{23} <$ {\color{red}{x.xxx}} ({\color{red}{x.xxx}}),
{\color{red}{x.xxx}} ({\color{red}{x.xxx}}) $< sin^{2}\theta_{13} <$ {\color{red}{x.xxx}} ({\color{red}{x.xxx}}) and
-$\pi$ (-$\pi$) $< \delta_{CP} <$ $\pi$ ($\pi$).

\end{abstract}

%
% Keywords
%

\end{frontmatter}

\clearpage

\tableofcontents

\clearpage

\listoffigures
\listoftables


